\uuid{IjiF}
\titre{Limite en un point}
\auteur{bodin ; hanani ; mzari}
\organisation{exo7}
\licence{CC BY-NC-SA 4.0}
\datecreate{2020-01-20}
\niveau{L1}
\module{Analyse}
\chapitre{Continuité, limite et étude de fonctions réelles}
\sousChapitre{Limite de fonctions}
\difficulte{2}
\theme{Limites des fonctions réelles, limite en un point}
\exo7qcmid{lille-1.limites.018}

\contenu{
\begin{question}
\qclassification{123.03}
\qtags{motcle=limite en un point}
Soit $f(x)= \frac{\sin x}{x}$. Quelles sont les assertions vraies ?
\begin{answers}

    \bad{$f$ n'admet pas de limite en $0$.}

    \good{$\lim_{x\to 0} f(x)=1$}

    \bad{$\lim_{x\to 0} f(x)=0$}

    \bad{$\lim_{x\to 0} f(x)=+\infty$}

\end{answers}
\begin{explanations}
La fonction $g: x\to \sin x$ est dérivable sur $\Rr$  et $g'(x)= \cos x$, pour tout $x\in \Rr$. Donc  $\lim_{x\to 0}f(x)= \lim_{x\to 0}\frac{g(x)-g(0)}{x-0} = g'(0)=1$.
\end{explanations}
\end{question}
}
